\section*{Capítulo \ref{ch:g8:the-oceans} --- Os oceanos}

\begin{solution}{exo:g8:the-oceans:1}
\omterm{def:g8:the-oceans:seafloor}{Plataforma continental}, \omterm{def:g8:the-oceans:seafloor}{talude continental}, \omterm{def:g8:the-oceans:seafloor}{planície abissal}, \omterm{def:g8:the-oceans:seafloor}{dorsal oceânica}.
\end{solution}

\begin{solution}{exo:g8:the-oceans:2}
(a) Falso: cerca de \qty{3700}{m}, ou seja, \qty{3.7}{km}. (b) Verdadeiro.
(c) Falso: a evaporação tira água e deixa os sais, então a água fica mais
salgada. (d) Verdadeiro.
\end{solution}

\begin{solution}{exo:g8:the-oceans:3}
(a) Uma \omterm{def:g8:the-oceans:seafloor}{dorsal oceânica}. (b) A \omterm{def:g8:the-oceans:salinity}{salinidade}. (c) Um \omterm{def:g8:the-oceans:current}{giro}. (d) Uma \omterm{def:g8:the-oceans:seafloor}{fossa oceânica}.
\end{solution}

\begin{solution}{exo:g8:the-oceans:4}
(a) No sentido anti-horário. (b) Quente: a Corrente do Golfo; fria: a
Corrente das Canárias. (c) Os \omterm{def:g8:winds-global-circulation:global-winds}{alísios}.
\end{solution}

\begin{solution}{exo:g8:the-oceans:5}
(a) $35 \div 1000 \times 100 = 3.5~\%$. (b) $20 \times 35 = \qty{700}{g}$.
(c) $1000 \div 35 \approx 28.6$: cerca de \qty{29}{kg}.
\end{solution}

\begin{solution}{exo:g8:the-oceans:6}
(a) $1500 \times 6 = \qty{9000}{m}$. (b) O som desce até o fundo e volta
para cima: ele percorre a profundidade duas vezes. (c) $9000 \div 2 =
\qty{4500}{m}$. (d) Uma \omterm{def:g8:the-oceans:seafloor}{planície abissal} (uma fossa seria mais funda, uma
dorsal ou uma plataforma, mais rasas).
\end{solution}

\begin{solution}{exo:g8:the-oceans:7}
(a) Entre cerca de \qty{3500}{m} e \qty{5000}{m}. (b) A cerca de \qty{4400}{km} ao longo do perfil, a cerca de \qty{2300}{m} de profundidade. (c) Cerca de
$4500 - 2300 = \qty{2200}{m}$ (entre 1200 e \qty{2700}{m}, conforme a
\omterm{def:g1:landscapes-around-us:landform}{planície} escolhida). (d) O \omterm{def:g2:water-on-earth:ocean}{oceano} tem alguns quilômetros de profundidade,
mas milhares de quilômetros de largura: desenhado na mesma escala, o
perfil inteiro seria mais fino que um traço.
\end{solution}

\begin{solution}{exo:g8:the-oceans:8}
O \omterm{def:g2:water-on-earth:ocean}{mar} A é mais salgado: o Sol e o ar seco causam uma forte
\emph{evaporação}, a água vai embora, \emph{os sais ficam}, e quase
nenhuma \omterm{def:g4:clouds-rain-snow:precipitation}{chuva} ou \omterm{def:g1:landscapes-around-us:water}{rio} traz \emph{\omterm{def:g2:water-on-earth:salt-fresh}{água doce}} de volta. O \omterm{def:g2:water-on-earth:ocean}{mar} B é menos
salgado: os seus \omterm{def:g1:landscapes-around-us:water}{rios} e a sua \omterm{def:g4:clouds-rain-snow:precipitation}{chuva} trazem muita \emph{\omterm{def:g2:water-on-earth:salt-fresh}{água doce}}, que
dilui os sais.
\end{solution}

\begin{solution}{exo:g8:the-oceans:9}
(a) A Corrente do Golfo, para o norte ao longo do \omterm{def:g1:landscapes-around-us:water}{litoral}, e depois para o
nordeste. (b) Para o leste (nordeste), com a Corrente do Atlântico Norte.
(c) Ao longo do \omterm{def:g1:landscapes-around-us:water}{litoral} do noroeste da África (e da Espanha e de
Portugal), com a Corrente das Canárias. (d) Fria.
\end{solution}

\begin{solution}{exo:g8:the-oceans:10}
(a) $6.4 \times 24 = 153.6$: cerca de \qty{154}{km}. (b) $153.6 \times 7
\approx \qty{1075}{km}$. (c) A Corrente do Golfo é um pouco mais rápida
que uma pessoa caminhando, mas ela nunca para, nem de dia nem de noite.
\end{solution}

\begin{solution}{exo:g8:the-oceans:11}
(a) \qty{3.7}{km} e \qty{11}{km}. (b) $11000 \div 3700 \approx 3.0$
vezes. (c) O som percorre $2 \times 11000 = \qty{22000}{m}$; $22000
\div 1500 \approx 14.7$: cerca de \qty{15}{s}. (d) $1000 = 10^3$ anos.
\end{solution}

\begin{solution}{exo:g8:the-oceans:12}
(a) Perto do \omterm{def:g3:continents-and-oceans:equator}{equador}, os \omterm{def:g8:winds-global-circulation:global-winds}{alísios} empurram a água tropical quente para o
oeste. Na borda oeste do Atlântico, ela vira para o norte com a Corrente
do Golfo, uma corrente quente, ao longo da Flórida e da América do Norte.
Os \omterm{def:g8:winds-global-circulation:global-winds}{ventos de oeste} a empurram então através do Atlântico em direção ao
nordeste, até os \omterm{def:g2:water-on-earth:ocean}{mares} ao largo da Noruega. (b) Ali, ela cede o seu calor
ao ar frio: ela esfria, e fica mais salgada quando se forma o gelo
marinho, porque o sal fica na água. A água fria e salgada é densa, então
ela afunda até o fundo. (c) Ela rasteja para o sul ao longo do fundo do
\omterm{def:g2:water-on-earth:ocean}{oceano}, em direção à Antártida, e depois entra nos \omterm{def:g2:water-on-earth:ocean}{oceanos} Índico e
Pacífico, onde sobe devagar de novo.
\end{solution}

\begin{solution}{pb:g8:the-oceans:1}
\textbf{1.} No \omterm{def:g2:water-on-earth:ocean}{oceano} Atlântico, no \omterm{def:g3:continents-and-oceans:equator}{hemisfério norte}.

\textbf{2.} No sentido horário; a boia é levada para o norte no começo.

\textbf{3.} O trecho A (Corrente do Golfo) é quente: ele vem dos
\omterm{def:g5:weather-and-climate:tropics}{trópicos}. O trecho C (Corrente das Canárias) é frio: ele vem de \omterm{def:g6:latitude-longitude:latitude}{latitudes}
mais altas, em direção ao \omterm{def:g3:continents-and-oceans:equator}{equador}.

\textbf{4.} Trecho B: os \omterm{def:g8:winds-global-circulation:global-winds}{ventos de oeste}. Trecho D: os \omterm{def:g8:winds-global-circulation:global-winds}{alísios}.

\textbf{5.} $3800 \div 100 = 38$ dias.

\textbf{6.} B: $2100 \div 20 = 105$ dias. C: $3100 \div 20 = 155$ dias.
D: $6100 \div 20 = 305$ dias.

\textbf{7.} $3800 + 2100 + 3100 + 6100 = \qty{15100}{km}$.

\textbf{8.} A Corrente do Golfo é a corrente estreita e rápida do lado
oeste do \omterm{def:g8:the-oceans:current}{giro}; as correntes dos outros lados são largas e lentas.

\textbf{9.} $38 + 105 + 155 + 305 = 603$ dias.

\textbf{10.} $603 \div 30 = 20.1$: cerca de \textbf{20 meses}.

\textbf{11.} Depois dos trechos A, B e C, já se passaram $38 + 105 + 155 = 298$
dias; os 365 dias chegam durante o trecho D (a Corrente Norte Equatorial), $365 -
298 = 67$ dias depois do seu começo.

\textbf{12.} $603 - 365 = 238$ dias a mais, ou seja, $238 \div 30 \approx 8$
meses a mais.
\end{solution}
